\section{Trabajo relacionado}
\label{sec:related-work}

\subsection{An\'alisis din\'amico basado en sandbox}

Los sandboxes automatizados de malware se remontan a Cuckoo
Sandbox~\cite{cuckoo}, cuya arquitectura (un controlador que orquesta
la detonaci\'on dentro de una m\'aquina virtual desechable, con un
agente en el hu\'esped que reporta llamadas a API y artefactos de
vuelta al host) sigue siendo el patr\'on est\'andar. CAPE~\cite{cape}
(Config And Payload Extraction) bifurc\'o este linaje espec\'ificamente
para a\~nadir desempaquetado autom\'atico y extracci\'on de
configuraci\'on para familias de malware conocidas, e incluye una
amplia biblioteca de firmas de comportamiento mantenida por la
comunidad, cada una opcionalmente etiquetada con identificadores de
t\'ecnica MITRE ATT\&CK. Esta pipeline se apoya directamente en CAPE
para el an\'alisis din\'amico (\S\ref{sec:dynamic}), y trata sus
mapeos de firma a t\'ecnica mantenidos por la comunidad como una
hip\'otesis de partida a verificar, no como datos de referencia a
asumir; una distinci\'on sustanciada de forma concreta en
\S\ref{sec:evaluation}: varios de los mapeos propios de CAPE se
comprobaron frente a su evidencia en bruto y resultaron ser
sistem\'aticamente demasiado amplios para las muestras en las que se
activaron, un paso de inspecci\'on que esta pipeline realiza y que la
mayor\'ia de despliegues de sandbox omiten.

\subsection{Mapeo est\'atico a ATT\&CK}

Mapear artefactos est\'aticos (importaciones, cadenas, coincidencias
YARA) directamente a t\'ecnicas ATT\&CK es una pr\'actica com\'un en
productos EDR comerciales y en herramientas de la comunidad. El
equivalente de c\'odigo abierto m\'as cercano es capa, de
Mandiant~\cite{capa}, que compara reglas curadas de importaciones,
cadenas y patrones de bytes contra un binario para reportar las
capacidades que implementa, muchas de ellas ya mapeadas a ATT\&CK. La
mayor\'ia de reglas de capa, como la mayor\'ia de b\'usquedas de
indicador de un EDR, son coincidencias de un \'unico indicador: basta
un nombre de importaci\'on o un patr\'on de cadena para activar la
regla. Este proyecto diverge deliberadamente de ese patr\'on en su
modelo de confianza (\S\ref{sec:static-rules}): un \'unico indicador
gen\'erico (una importaci\'on usada para muchos prop\'ositos
leg\'itimos) se trata como evidencia m\'as d\'ebil que una
\emph{combinaci\'on coherente} de indicadores que solo tiene sentido en
conjunto para la t\'ecnica afirmada. Concretamente,
\technique{OpenProcess} por s\'i sola no es evidencia de inyecci\'on de
proceso (\technique{T1055}), pero \technique{OpenProcess} junto con
\technique{VirtualAllocEx} y \technique{WriteProcessMemory} s\'i lo es,
porque las tres API componen una primitiva reconocible que no tiene
otro prop\'osito com\'un. Esto sacrifica parte de la amplitud de capa
(que incluye cientos de reglas cubriendo muchas plataformas y clases de
capacidad) a cambio de una precisi\'on m\'as estrecha pero respaldada
por evidencia en las t\'ecnicas de Windows que esta pipeline s\'i
cubre, una decisi\'on de compromiso explicitada y medida directamente
en \S\ref{sec:fp-audit}.

\subsection{Extracci\'on autom\'atica desde informes de amenazas}

Una l\'inea de trabajo paralela extrae t\'ecnicas ATT\&CK no de un
binario sino de prosa de analista sin estructurar: TTPDrill~\cite{ttpdrill}
analiza informes de inteligencia de amenazas con una ontolog\'ia
propia y una pipeline de PLN para recuperar acciones del atacante;
TRAM, de MITRE Engenuity~\cite{tram}, y rcATT~\cite{rcatt} entrenan
ambos clasificadores de texto para etiquetar p\'arrafos de informes CTI
con etiquetas de t\'ecnica ATT\&CK probables. Estas herramientas
resuelven un problema genuinamente distinto: ayudan a un analista que
ya dispone de un \emph{informe escrito} a convertirlo en etiquetas de
t\'ecnica estructuradas y legibles por m\'aquina de forma m\'as
r\'apida. Esta pipeline, en cambio, deriva sus etiquetas de t\'ecnica
directamente del propio binario y comportamiento en tiempo de
ejecuci\'on de la muestra, sin ning\'un informe como paso intermedio,
lo que permite que cada hallazgo lleve asociada una pieza concreta de
evidencia verificable (una combinaci\'on de importaciones, el campo de
datos en bruto de una firma) en lugar de la puntuaci\'on de confianza
de un clasificador de PLN sobre texto libre. Los dos enfoques son
complementarios, no competidores: un clasificador al estilo TRAM
podr\'ia en principio tomar como texto de entrada los propios casos de
estudio escritos por esta pipeline (\S\ref{sec:case-studies}), cerrando
el ciclo de evidencia binaria a informe y de vuelta a etiquetas de
t\'ecnica reextra\'idas.

\subsection{Agregadores de reputaci\'on y esc\'aneres multimotor}

VirusTotal~\cite{virustotal} es la herramienta a la que primero recurre
la mayor\'ia de analistas de respuesta a incidentes, y constituye un
punto de comparaci\'on \'util porque opera en el extremo opuesto del
compromiso entre amplitud y profundidad que asume esta pipeline.
VirusTotal agrega veredictos de decenas de motores antivirus e
integraciones de sandbox sobre un corpus de env\'ios de una escala que
ning\'una pipeline de investigaci\'on individual puede igualar, lo que lo
hace excelente para el triaje r\'apido de reputaci\'on: si un hash ya es
conocido, y si otros proveedores ya lo marcan. Lo que VirusTotal no
hace es reconciliar esas se\~nales entre s\'i para una \'unica
afirmaci\'on: el veredicto de cada motor y las etiquetas de
comportamiento de cada sandbox se muestran uno junto a otro, no se
contrastan, de modo que dos etiquetas de t\'ecnica contradictorias
procedentes de dos integraciones distintas se muestran ambas sin m\'as,
sin ning\'un mecanismo que resuelva cu\'al de ellas respalda realmente
la evidencia en bruto. Esta pipeline apunta precisamente a ese hueco,
no a la amplitud de VirusTotal: muchas menos muestras, pero cada
afirmaci\'on de t\'ecnica trazada hasta la combinaci\'on de
importaciones, el campo de datos de firma o el patr\'on de cadena
exacto que la produjo, la concordancia entre fuentes y entre niveles
modelada expl\'icitamente (\S\ref{sec:reconciliation}), y los
componentes soltados por una muestra multietapa analizados y atribuidos
de forma individual en lugar de fundirse en un \'unico veredicto
agregado para el hash del padre (\S\ref{sec:resubmission}). Las dos
herramientas no son sustitutas: una consulta a VirusTotal es el primer
paso correcto para el triaje a escala, y esta pipeline es el siguiente
paso correcto una vez que una muestra concreta necesita un mapeo de
t\'ecnica defendible y ligado a evidencia, no una puntuaci\'on de
reputaci\'on.

\subsection{Est\'andares de intercambio de inteligencia de amenazas}

STIX~2.1~\cite{stix21} (Structured Threat Information Expression) y
MISP~\cite{misp} (Malware Information Sharing Platform) son los dos
formatos de intercambio a los que exporta esta pipeline. STIX define
un grafo de objetos tipados (indicadores, malware, patrones de ataque,
relaciones) pensado para el intercambio m\'aquina a m\'aquina; MISP es a
la vez un modelo de datos y una plataforma en ejecuci\'on que la
mayor\'ia de organizaciones ya operan. Exportar a ambos, en lugar de
producir \'unicamente un informe legible por humanos, es lo que hace
que el resultado de la pipeline sea directamente consumible por las
herramientas que un equipo de operaciones de seguridad ya utiliza, sin
necesidad de transcripci\'on manual.

\subsection{Posicionamiento}

Esta tesis no propone una primitiva de detecci\'on novedosa: cada
se\~nal individual utilizada (una importaci\'on de API, un patr\'on de
cadena, una firma de comportamiento de sandbox) es una t\'ecnica bien
establecida en la pr\'actica del an\'alisis de malware, y herramientas
como capa y CAPE ya exponen muchas de esas mismas se\~nales en bruto.
La contribuci\'on est\'a en la \emph{integraci\'on}: un modelo de
reconciliaci\'on de confianza que es verificablemente m\'as que la suma
de sus fuentes (\S\ref{sec:reconciliation}), un mecanismo para atribuir
correctamente la capacidad distribuida del malware multietapa
(\S\ref{sec:resubmission}), y, argumentada como la m\'as defendible de
las tres, una disciplina de validaci\'on en la que cada afirmaci\'on que
hace la pipeline se comprueba frente a evidencia real antes de confiar
en ella, documentada con suficiente especificidad como para que la
comprobaci\'on misma sea reproducible (\S\ref{sec:evaluation}).
Posicionada frente a los dos polos revisados arriba, esta pipeline se
sit\'ua deliberadamente entre un agregador amplio y no reconciliado
como VirusTotal y un informe manual de ingenier\'ia inversa estrecho,
de una sola muestra: automatizado y exportador de est\'andares como el
primero, verificado por evidencia y espec\'ifico por t\'ecnica como el
segundo.
